
This paper documents a failure mode in an automated research pipeline: an experiment that was fully implemented but could not execute due to computational constraints not validated during the design phase. The specific case involved investigating LoRA-MoE coordination mechanisms for multi-task learning using Mixtral-8x7B (47 billion parameters). The pipeline completed all implementation phases—generating 29 Python files, 10 test suites, passing all code quality checks—before attempting execution and encountering CUDA Out of Memory errors. Analysis revealed the model required approximately 489GB VRAM while available capacity was 475GB across five H100 GPUs.

This failure, while specific to one automated system, reveals a structural gap: the YouRA pipeline contained no systematic checkpoint between experiment design (Phase 2C) and implementation (Phase 3-4) to validate computational feasibility. Experiment designs specifying model and dataset choices proceeded directly to implementation planning and coding without verifying resource requirements against available hardware. Traditional software quality metrics (test passing, task completion, code coverage) detected no issues because computational feasibility is orthogonal to implementation correctness.

The gap appears systematic to this pipeline architecture. Phase 2C focuses on scientific requirements—model selection prioritized theoretical needs (native MoE architecture, sufficient model capacity for the coordination hypothesis) without explicit hardware constraint validation. Phase 3-4 assumes approved designs are executable—if Phase 2C outputs a specification, implementation proceeds. No checkpoint asks: "Will this configuration fit available hardware before we invest implementation effort?"

We quantify the cost in this case: 10-16 hours of implementation effort (Phase 3: 2-3 hours planning, Phase 4: 8-13 hours coding and testing) occurred before discovering infeasibility. We propose a Phase 2C.5 computational feasibility gate that estimates memory requirements using a formula accounting for major components (model parameters, optimizer states, gradients, activations, framework overhead) and validates against available hardware capacity. The estimated overhead is 5 minutes. Retrospective application to this failure case shows the gate would have correctly predicted infeasibility (489GB required vs 475GB available, 103\% utilization exceeding 85\% threshold).

This contribution is process-level rather than algorithmic. We identify a missing validation step in one automated research workflow, document one failure case with quantified cost, and propose a solution design with retrospective validation on that case. Generalizability across diverse workflows and empirical validation of gate effectiveness requires multi-project deployment.

The paper proceeds as follows: Section 2 reviews related work in model optimization, workflow tools, and negative results reporting. Section 3 describes the proposed feasibility gate design with memory estimation formula. Section 4 presents the experimental context (what was attempted and why it failed). Section 5 provides evidence of the workflow gap through implementation metrics and cost analysis. Section 6 discusses limitations, implications, and future work.
